% !TeX program = pdflatex
\documentclass[sigconf,anonymous,review]{acmart}
\usepackage{CJKutf8}
\usepackage{amsmath}
\usepackage{booktabs}
\setlength{\emergencystretch}{2em}
\sloppy
\settopmatter{printfolios=true}
\setcopyright{none}
\renewcommand{\shortauthors}{Anonymous Authors}
\title[GSU]{GSU：关联结构引导的多跳知识遗忘}
\author{Anonymous Authors}
\begin{document}
\begin{CJK*}{UTF8}{gbsn}
\begin{abstract}
对原始问答的遗忘不等于目标知识无法通过关联查询被恢复。本文研究关联结构引导的知识遗忘算法 GSU：通过模型探测与贪心覆盖识别候选支持事实，将关联信息转化为组件和更新方向指导，检验仅对原始事实施加遗忘监督能否抑制多跳回忆。本文建立直接、改写和关联多跳评测，并通过结构化引导与随机对照分析更新机制。当前单随机种子的筛选结果表明，直接加入支持事实仍比已有的仅原始事实组件引导更有效；软引导未达到预期，硬引导呈现值得进一步检验的差异。本文草稿据此明确后续方法验证与证据边界，尚不声称 GSU 已实现仅原始事实监督下的可靠多跳遗忘。
\end{abstract}
\keywords{知识遗忘，多跳回忆，关联结构，梯度指导，参数高效微调}
\maketitle

\section{引言}
网络文本中的事实通过实体、关系和支持知识相互连接。模型对某条事实的直接回答被削弱后，仍可能回答涉及关联结构的多跳问题。因此，面向知识撤回的遗忘评估需要同时考察直接访问和间接回忆，而不能仅检查训练问题的正确率。

直接扩大遗忘数据是一种可行路径，但它同时改变了原始目标之外的支持事实。本文关注另一问题：能否先识别支持目标回忆的关联结构，再将这些信息嵌入遗忘更新，使遗忘监督仍来自原始目标事实？该问题将关联知识作为算法指导，与将其作为直接遗忘对象区分开来。

本文将事实扩展、组件定位和梯度更新作为相互关联而又需要分别验证的环节。事实扩展提供候选支持信息；组件定位确定更新位置；方向干预决定如何写入更新。任何一个代理信号的改善，都需要通过实际参数位移、答案概率与最终多跳行为进行核验。相关工作的正式文献比较将在后续补充可核验引用，本稿不主张优先性。

本文以持续优化的 GSU 为算法载体，组织关联探测、更新指导和行为诊断。拟验证的贡献包括：原始监督与关联指导分离的问题设定；关联结构引导的更新机制；以及直接、多跳和影响范围的分层验证。当前研究仍处于机制筛选阶段，正文分别标注已有观察与待验证设计，避免把预期贡献写成完成结果。

\section{问题定义}
设冻结模型为 $M_{\theta_0}$，遗忘后模型为 $M_\theta$。$F_0$ 是原始待遗忘事实，$F_1$ 是通过探测与贪心覆盖选出的关联支持事实；$Q_{\mathrm{multi}}$ 是多跳评测查询，而非 $F_1$ 本身。目标方法仅对 $F_0$ 施加遗忘损失，允许使用预先声明的关联事实作为定位信息，但不以最终多跳评测题作为梯度目标。

本文以直接问题、留出改写和多跳问题的回忆率衡量目标可访问性。三选一题目的低正确率不等于知识已从参数中彻底删除；答案概率和候选 margin 作为辅助证据。保留知识按目标关联程度进行审计，不能从总保留率直接判断所有下降均为误伤或必要代价。

\section{GSU 方法框架}
\subsection{关联结构识别}
从原始事实簇的图路径中抽取候选支持事实，以模型探测估计掌握情况，并依据已有多跳探测记录标记活跃路径。在固定事实预算下，以贪心覆盖选择 $F_1$。既有选择使用了路径探测元数据，因此当前覆盖链测试属于机制诊断，正式泛化结论还需要独立查询和路径划分。

\subsection{组件与梯度指导}
当前已完成的候选在 LoRA 参数空间中计算组件分数。每个组件对应某层目标矩阵上的 LoRA A/B 参数对。令 $g_{ij}$ 为第 $i$ 条原始事实在组件 $j$ 上的答案对数概率梯度，$\bar g_j$ 为其均值；令 $h_{mj}$ 为关联事实的对应梯度，$c_m$ 为固定活跃路径数。定义
\begin{align}
 A_j &= \frac{\|\sum_i g_{ij}\|}{\sum_i\|g_{ij}\|},\\
 P_j &= \frac{\sum_m c_m\cos(\bar g_j,h_{mj})}{\sum_m c_m},\\
 S_j &= A_j\max(P_j,0).
\end{align}
零范数时相应项取零。软指导以 $S_j/\max_k S_k$ 重加权梯度，硬指导保留不低于分数中位数的组件；随后恢复整体梯度范数。分数仅在训练前计算并冻结。全零软分数退化为恒等权重。这一分数是已测试构造，不代表最终 GSU 必须采用它。

\subsection{实际更新与反馈}
组件梯度经过 AdamW 后的实际位移不必与原始梯度同向，因此记录原始更新和反馈后状态。当前 FB16 控制器测量保留锚点的实际答案概率下降，并在固定预算下修正；它不是直接监测多跳目标的反馈控制器。面向目标回忆的最终反馈机制仍待设计与独立验证，不能把当前保留反馈写成已经实现的多跳行为闭环。

\subsection{参数化与机制边界}
当前 1.5B 模型有 28 层，每层七类目标矩阵，共 196 个组件。LoRA 的低秩参数化、初始化和优化器状态均可能影响梯度分数与实际更新。硬掩码只作用于遗忘梯度，不能解释为相关参数绝对冻结；权重衰减和后续反馈仍可能修改它们。

\section{实验}
\subsection{设置与数据边界}
已有机制筛选使用 Qwen2.5-1.5B-Instruct、FaithUn 固定 split、seed 42、LoRA 秩 8、缩放参数 16、学习率 $5\times10^{-5}$，各组统一 114 次更新。$F_0$ 包含 33 条直接事实，$F_1$ 包含五条支持事实。每步从前 33 条 retaining 中抽取 16 条锚点；反馈最多两轮，修正预算为原始位移范数的 0.30。

官方 forget、retaining、test 分别包含 33、66、464 个事实簇，索引互斥不意味着实体或路径独立。前 33 条 retaining 用于保护，后 33 条用于未参与锚点保护的审计；后者也不是已证明语义独立的集合。

\subsection{RQ1：直接遗忘是否留下多跳回忆？}
比较仅原始监督和扩展监督。表~\ref{tab:screen} 中 A0 的关联多跳命中为 52/75，A1 为 22/75，说明在本次配置下扩展监督进一步降低了该切片的回忆率。该观察不能独自证明所有答题都通过所选路径推导。

\begin{table}[t]
\caption{已完成的单种子筛选。多跳与直接命中越低表示目标回忆越弱。D 为非 F1 重叠诊断切片，越高表示该切片的保持越好。}
\label{tab:screen}
\centering\small
\begin{tabular}{lrrrr}
\toprule
方法 & 多跳/75 & F0/33 & D/12 & retain/33\\
\midrule
A0：F0 & 52 & 8 & 9 & 24\\
A1：F0+F1 & 22 & 7 & 3 & 19\\
A2：F0 软指导 & 57 & 9 & 9 & 18\\
A3：F0 硬指导 & 47 & 6 & 1 & 19\\
A4：扩展软指导 & 21 & 3 & 2 & 16\\
A2R：随机软指导 & 64 & 10 & 9 & 23\\
A3R：随机硬指导 & 73 & 13 & 9 & 28\\
\bottomrule
\end{tabular}
\end{table}

\subsection{RQ2：仅 F0 的结构指导是否有效？}
A2 的多跳命中高于 A0，未显示预期改善。A3 比 A0 少命中五题，但仍远高于扩展监督的 A1。A4 比 A1 少命中一题，不能据此确认软指导的稳定增益。已有结果尚不支持当前组件指导替代 $F_1$ 直接监督。

\subsection{RQ3：结构化指导是否优于随机？}
A3 与随机硬指导 A3R 分别命中 47/75 与 73/75，提供了需重复验证的差异。随机组同时有更弱的 F0 遗忘，因而该差异可能混合整体遗忘强度效应；后续需要在匹配 F0 遗忘程度的条件下比较，不能直接归因为多跳路径定位成功。

\subsection{RQ4：哪些更新机制起作用？}
拟对最终方法的关联探测、组件先验、方向干预和反馈分别消融，并记录梯度到实际位移再到行为变化的链条。历史组件撤销用于提供机制候选，不直接作为新方法的成功证据。既有探索版本按必要性放入附录，正文不按时间顺序堆叠候选。

\subsection{RQ5：模型规模与稳定性如何？}
拟使用共同题目、统一 $F_0/F_1$ 和多个训练种子对 1.5B/3B 配对评测，并按跳数、路径密度和桥接实体类型分层。已有 clean 结果中原始 hop2 约为 0.644/0.633，不能预设 3B 多跳能力更强。跨模型关联事实选择和样本掌握程度必须单独控制。

\subsection{RQ6：影响是否集中于目标关联知识？}
拟将审计知识按同实体、共享答案、路径邻近等标签分层，再比较低关联样本。结构重叠只能提出关联线索，不能证明模型实际调用该知识；需要匹配初始难度与关系类型，并结合干预验证。影响范围审计用于排除广泛能力退化，不替代目标回忆的主要评测。

\section{相关工作（待补引用）}
正式稿将比较事实级知识遗忘、多跳回忆评估、关联事实筛选、组件定位和梯度方向控制。对比须明确原始方法与本地 LoRA 移植的差别、监督数据权限和计算预算；本稿未插入未经核验的参考文献条目。

\section{讨论与结论}
本文将关联结构作为遗忘算法的指导信息，围绕 F0-only 监督下的多跳回忆抑制组织 GSU。现有筛选显示直接扩展监督仍是主要有效因素，当前组件指导尚不足以替代它。最终论文需要补齐匹配随机对照、多种子、跨模型以及独立查询上的证据。本文的“遗忘”指实际测试条件下的回忆抑制，不代表所有潜在恢复方式均被消除。

\appendix
\section{草稿状态与结果溯源}
本文件为中文内部工作稿，使用 CJKutf8 和 pdfLaTeX；不作为正式英文投稿文件。表中数字来自 exp006 批次 \nolinkurl{batch_20261021_132023_431bf0} 的各组 \nolinkurl{metrics.json} 和 \nolinkurl{comparison.json}。已有观察与拟实施实验在正文中分别说明；后续修改不得覆盖历史运行记录。
\end{CJK*}
\end{document}
